% =====================================================================
% Methods section — vertical velocity assimilation on Thwaites
% Style follows Shapero et al. (2021), GMD, "icepack: a new glacier
% flow modeling package in Python, version 1.0".
% =====================================================================

\section{Methods}
\label{sec:methods}

We estimate the depth-averaged ice fluidity $A(x,y)$ and basal
friction coefficient $C(x,y)$ on Thwaites Glacier by minimizing a
single objective functional that combines the misfit to surface
velocity observations from satellite radar interferometry with the
misfit to vertical strain rates measured at a network of ApRES
(autonomous phase-sensitive radio echo sounder) sites. The forward
model is the hybrid (Blatter--Pattyn) flow model implemented in
icepack \citep{Shapero2021}, discretized on an extruded mesh with a
spectral basis in the vertical. The inverse problem is solved by
quasi-Newton minimization with adjoint-derived gradients computed
through the algorithmic differentiation framework
\texttt{tlm\_adjoint} \citep{Maddison2019}.

% ---------------------------------------------------------------------
\subsection{Forward model}
\label{sec:forward}

The hybrid model solves for the horizontal velocity
$\bm u(x,y,\zeta) = (u, v)$ on the terrain-following coordinate
$\zeta = (z-b)/h$, where $b$ is the bed elevation and
$h = s - b$ the ice thickness. Following \citet{Shapero2021}, the
velocity is a critical point of the action functional
\begin{equation}
\label{eq:action}
J(\bm u) = \mathrm{viscosity} + \mathrm{friction}
        - \mathrm{gravity} - \mathrm{terminus},
\end{equation}
in which
\begin{align}
\mathrm{viscosity} &=
   \frac{n}{n+1}\!\iint_\Omega\!\!\int_0^1
   h\, A^{-1/n}\!\left(|\dot{\bm\varepsilon}_x(\bm u)|^2
                       + h^{-2}|\partial_\zeta \bm u|^2\right)^{(1+1/n)/2}
   \mathrm d\zeta\,\mathrm d x,
   \label{eq:visc}\\[2pt]
\mathrm{friction} &=
   \frac{m}{m+1}\!\int_\Omega C\,\varphi\,
                  |\bm u(\zeta\!=\!0)|^{1+1/m}\,\mathrm d x,
   \label{eq:fric}\\[2pt]
\mathrm{gravity} &=
   -\!\iint_\Omega\!\!\int_0^1 \rho_\mathrm I g\,h\,
                  \nabla s\!\cdot\!\bm u\,\mathrm d\zeta\,\mathrm d x,
   \label{eq:grav}\\[2pt]
\mathrm{terminus} &=
   \!\iint_\Gamma\!\!\int_0^1 (p_\mathrm I - p_\mathrm W)\,
                 \bm u\!\cdot\!\bm\nu\,\mathrm d\zeta\,\mathrm d\gamma,
   \label{eq:term}
\end{align}
where $n=3$ is the Glen flow law exponent, $m=3$ is the Weertman
sliding exponent, $\rho_\mathrm I$ and $\rho_\mathrm W$ are ice and
sea-water densities, $g$ is gravity, and $\bm\nu$ is the outward
normal at the calving front $\Gamma$. The horizontal strain-rate
tensor in terrain-following coordinates is
\begin{equation}
\dot{\bm\varepsilon}_x(\bm u)
   = \tfrac{1}{2}\left(\nabla_\zeta\bm u
                       + (\nabla_\zeta\bm u)^{\!\top}\right)
     + \tfrac{1}{2}\left(\partial_\zeta\bm u\otimes\bm v_\zeta
                         + \bm v_\zeta\otimes\partial_\zeta\bm u\right),
   \label{eq:eh}
\end{equation}
with the geometric correction vector
$\bm v_\zeta = -h^{-1}[(1-\zeta)\nabla b + \zeta\nabla s]$
arising from the terrain-following chain rule
\citep[][Eq.~19]{Shapero2021}.

The factor $\varphi$ in the friction term~\eqref{eq:fric} is the
flotation fraction
\begin{equation}
\varphi(x,y) = \max\!\left(0,\,
   1 - \frac{\rho_\mathrm W\,g\,\max(0,\,h-s)}{\rho_\mathrm I\,g\,h}\right),
   \label{eq:phi}
\end{equation}
which enforces zero basal traction beneath floating ice. The water
back-pressure at the calving front in~\eqref{eq:term} follows the
piecewise-linear profile of \citet{Shapero2021},
$p_\mathrm W = \rho_\mathrm W g h\,(\zeta_\mathrm{sl} - \zeta)_+$.

% ---------------------------------------------------------------------
\subsection{Discretization}
\label{sec:discretization}

We mesh the Thwaites footprint $\Omega$ with triangular elements
generated by Gmsh \citep{Geuzaine2009}, using a target edge length
of \SI{2}{\km} that grades to \SI{10}{\km} away from the
fast-flowing trunk. The 2D triangulation $\mathcal T_h$ is then
extruded to a single-layer prismatic mesh on
$\Omega\times[0,1]$ using the \texttt{ExtrudedMesh} utility of
Firedrake \citep{Rathgeber2016}. We employ tensor-product elements:
continuous Lagrange polynomials of degree $h_d=1$ in the horizontal
combined with Gauss--Legendre polynomials of degree $v_d$ in the
vertical for the velocity space
\begin{equation}
\bm V_{h_d, v_d} = \mathrm{CG}_{h_d}(\Omega)\otimes
                    \mathrm{GL}_{v_d}([0,1]),\quad
\bm u_h \in \bm V_{h_d, v_d}^2.
   \label{eq:vspace}
\end{equation}
Setting $v_d=1$ gives a velocity that is linear in $\zeta$, while
$v_d=2$ resolves quadratic vertical structure. By incompressibility
the corresponding vertical-velocity field
$w(\zeta) = -\!\int_0^\zeta(\partial_x u + \partial_y v)\,h\,
\mathrm d\zeta'$ is one degree higher in $\zeta$ than the horizontal
velocity, so $v_d=1$ already accommodates a quadratic
$w$-profile.

Geometry fields ($h$, $s$, $b$) and the surface velocity
observations $\bm u_\mathrm{obs}$ are interpolated from BedMachine
Antarctica v4.1 \citep{Morlighem2020} and MEaSUREs
\SI{450}{\m} ice velocity v2 \citep{Mouginot2019} onto the 2D base
mesh, then lifted to the extruded mesh as $\zeta$-constant fields
through $\mathrm{CG}_1\otimes\mathrm{DG}_0$. The boundary
$\partial\Omega$ is partitioned into a calving segment $\Gamma$
(natural Neumann condition with hydrostatic ocean back-pressure) and
an inflow segment $\partial\Omega_\mathrm{in}$ where we impose a
Dirichlet condition $\bm u = \bm u_\mathrm{obs}$.

% ---------------------------------------------------------------------
\subsection{Observational data}
\label{sec:data}

\paragraph{Surface velocity.} The MEaSUREs annual mosaic supplies
the horizontal surface velocity $\bm u_\mathrm{obs}(x,y)$ and the
component-wise standard deviations $\sigma_x(x,y)$, $\sigma_y(x,y)$
on a \SI{450}{\m} grid. We retain only points whose reported
uncertainty is less than \SI{1e5}{\meter\per\year} and bound the
floor at \SI{1}{\meter\per\year} to avoid singular weighting. The
remaining $N_\mathrm{vel}$ scattered points are interpolated onto a
\texttt{VertexOnlyMesh} \citep{Nixon-Hill2024} so that the misfit can
be evaluated point-wise on the 2D base mesh.

\paragraph{ApRES vertical strain rates.} Repeat-pass ApRES surveys
acquired during the GHOST 2022--2023 field campaign provide vertical
velocity profiles $w(z)$ at $N_\mathrm{site}=91$ sites distributed
across the trunk and tributaries of Thwaites
\citep{Hoffman2024}. At each site we work in the local
terrain-following coordinate $\zeta = 1 - z/H$ where $H$ is the
ApRES-measured ice thickness, which can differ from the BedMachine
thickness by several hundred metres near the grounding line. The
uncertainty $\sigma_w(z)$ on each vertical-velocity sample is
propagated from the cross-correlation noise of the radar returns.
We retain samples with $z \in [z_\mathrm{min}, z_\mathrm{max}]$
that exclude the firn layer ($z_\mathrm{min}=\SI{100}{\m}$) and
remain within the reliable radar range; the choice of
$z_\mathrm{max}$ is discussed in
Sec.~\ref{sec:inverse}. Vertical strain rate samples are computed
by central differences of the vertical-velocity profile,
\begin{equation}
\dot\varepsilon_{zz}^\mathrm{obs}(z_i)
   = \frac{w(z_{i+1})-w(z_{i-1})}{z_{i+1}-z_{i-1}},
   \quad
\sigma_{\varepsilon}^2(z_i)
   = \frac{\sigma_w^2(z_{i+1})+\sigma_w^2(z_{i-1})}
          {(z_{i+1}-z_{i-1})^2}.
   \label{eq:apres-eps}
\end{equation}

% ---------------------------------------------------------------------
\subsection{Inverse problem}
\label{sec:inverse}

We seek the log-fluidity and log-friction fields
\begin{equation}
\theta_A(x,y) = \log(A/A_0),\qquad
\theta_C(x,y) = \log(C/C_0),
   \label{eq:controls}
\end{equation}
relative to background scales
$A_0 = \SI{20}{\mega\pascal^{-3}\per\year}$ and
$C_0 = \SI{e-2}{\mega\pascal\,(\meter\per\year)^{-1/m}}$. Both
controls are continuous Lagrange of degree $1$ on $\mathcal T_h$ and
are lifted to the extruded mesh as $\zeta$-constant fields. The
friction control is multiplied by the grounded mask
$\mathbb 1[\varphi > 0.01]$ to suppress its action on floating ice
during the optimization while leaving the basal-friction term
naturally zero through Eq.~\eqref{eq:phi}.

The total objective functional is
\begin{equation}
\mathcal L(\theta_A,\theta_C)
  = \mathcal J_\mathrm{vel}(\bm u)
  + \mathcal J_\mathrm{ApRES}(\bm u)
  + \mathcal R(\theta_A,\theta_C),
   \label{eq:total}
\end{equation}
where $\bm u$ solves the forward problem~\eqref{eq:action} for the
current $(\theta_A, \theta_C)$. The three components are described
below.

\paragraph{Surface velocity misfit.} On the extruded mesh the
surface horizontal velocity is recovered exactly by evaluating
$\bm u_h$ at $\zeta = 1$. We instead use the depth-integrated
$L^2$ form because it does not require sampling on the upper
boundary,
\begin{equation}
\mathcal J_\mathrm{vel}
  = \frac{1}{2|\Omega|}\!\iint_\Omega\!\!\int_0^1
       \frac{|\bm u_h - \bm u_\mathrm{obs}|^2}{\sigma_u^2}
       \,\mathrm d\zeta\,\mathrm d x,
   \label{eq:Jvel}
\end{equation}
with a uniform observational scale
$\sigma_u = \SI{10}{\meter\per\year}$ that broadly matches the
trunk-line uncertainty of MEaSUREs. The depth integral is exact
for the lifted observation field
($\bm u_\mathrm{obs}(x,y,\zeta) = \bm u_\mathrm{obs}(x,y)$) and
yields the depth-averaged horizontal velocity in the limit
$|\bm u(\zeta)|\to\bar{\bm u}(x,y)$.

\paragraph{ApRES strain-rate misfit.} The model vertical
strain rate is obtained from incompressibility,
\begin{equation}
\dot\varepsilon_{zz}^\mathrm{mod}(\bm u) =
   -\,\mathrm{tr}\,\dot{\bm\varepsilon}_x(\bm u),
   \label{eq:eps_zz}
\end{equation}
with $\dot{\bm\varepsilon}_x$ from Eq.~\eqref{eq:eh}. To compare
$\dot\varepsilon_{zz}^\mathrm{mod}$ with the ApRES samples we
construct a 3D \texttt{VertexOnlyMesh} on the extruded mesh whose
vertices are the points $(x_i, y_i, \zeta_i)$ obtained by stacking
each ApRES site over its retained depths. Letting $\Delta$ be
the resulting $\mathrm{DG}_0$ space on the vertex mesh, the
strain-rate misfit reads
\begin{equation}
\mathcal J_\mathrm{ApRES}
  = \frac{w_\mathrm{ApRES}}{2 N}
    \sum_{i=1}^{N}
    \frac{[\dot\varepsilon_{zz}^\mathrm{mod}(x_i,y_i,\zeta_i)
         - \dot\varepsilon_{zz}^\mathrm{obs}(x_i,y_i,\zeta_i)]^2}
         {\sigma_\varepsilon^2(x_i,y_i,\zeta_i)},
   \label{eq:Japres}
\end{equation}
where $N = N_\mathrm{site}\times N_z$ is the total number of
(site, depth) samples and $w_\mathrm{ApRES}$ is a tunable weight that
allows the relative influence of the ApRES term to be adjusted
without changing the data normalization. The interpolation
$\dot\varepsilon_{zz}^\mathrm{mod}\big|_{(x_i,y_i,\zeta_i)}$ is
performed by a Firedrake interpolation operator from the extruded
mesh into $\Delta$, which is differentiable with respect to the
controls through the adjoint of the forward solve.

The choice of $\zeta_i = 1 - z_i / H_\mathrm{ApRES}$ uses the
\emph{ApRES-measured} thickness rather than the BedMachine
thickness so that the relative depth coordinate is consistent
between the model evaluation and the data. We restrict the
analysis to $z\in[100,800]\,\mathrm m$, which excludes firn
compaction signals at the surface and stays within the depth range
where the radar return is reliable.

\paragraph{Regularization.} We use Tikhonov regularization on the
gradients of both controls,
\begin{equation}
\mathcal R(\theta_A,\theta_C)
  = \frac{L^2}{2|\Omega|}\!\int_\Omega
       \left(|\nabla\theta_A|^2 + |\nabla\theta_C|^2\right)\!
       \mathrm d x,
   \label{eq:reg}
\end{equation}
with a single correlation length scale $L$. We take
$L = \SI{10}{\km}$, comparable to the local mesh resolution; this
choice places the corner of the L-curve at a value of $L$ for which
the misfit terms~\eqref{eq:Jvel}--\eqref{eq:Japres} are not yet
limited by the prior.

% ---------------------------------------------------------------------
\subsection{Optimization}
\label{sec:optim}

The objective functional~\eqref{eq:total} is a function of the
finite-dimensional control vector
$\bm\theta = (\theta_A^h, \theta_C^h) \in \mathbb R^{2 N_\theta}$
with $N_\theta = \dim\,\mathrm{CG}_1(\mathcal T_h)$. We compute the
gradient $\nabla_{\!\bm\theta}\mathcal L$ by reverse-mode
algorithmic differentiation through the entire forward solve,
including the assembly of $\mathcal J_\mathrm{ApRES}$ via the
3D \texttt{VertexOnlyMesh} interpolation. The differentiation tape
is managed by \texttt{tlm\_adjoint} \citep{Maddison2019}; the
regularization gradient is added analytically by assembling
$L^2 |\Omega|^{-1} (\nabla\theta_\bullet,\,\nabla\!\cdot\,)$
against the test function of $\mathrm{CG}_1(\mathcal T_h)$.

The composite gradient is fed to the L-BFGS-B implementation of
\texttt{scipy.optimize.minimize} \citep{Virtanen2020} with a memory
of $m=10$, function tolerance $10^{-12}$, and gradient tolerance
$10^{-8}$. Each forward solve is performed by the
\texttt{FlowSolver} class of icepack with a damped Newton iteration
and a direct (MUMPS) factorization of the Jacobian
\citep{Amestoy2001}. Convergence failures of the forward Newton
iteration are caught by the optimization driver: in those cases the
objective is returned as a finite penalty and the previous
gradient is reused as a back-tracking direction so that L-BFGS
retreats along its line search rather than terminating.

For the high-resolution mesh we warm-start $\theta_A$ and
$\theta_C$ from a converged inversion that uses surface velocity
alone, which reduces the wall-clock time to convergence by
approximately one order of magnitude. The same warm start is used
for both vertical-degree experiments ($v_d=1$ and $v_d=2$),
allowing a direct comparison of their ability to reproduce the
ApRES observations under identical control initializations.
